
\begin{table}[htbp]
    \centering
    \caption{Identity and Ontological Determination}
    \label{tab:identity-and-determination}

    \renewcommand{\arraystretch}{1.4}
    \begin{tabular}{p{3.2cm} p{5.1cm} p{5.3cm}}
        \hline
        \textbf{Dimension} &
        \textbf{Ontological Characterization} &
        \textbf{Implication for Identity} \\
        \hline

        Being and determination &
        Being is not an indeterminate substrate that first exists and subsequently receives determination. To be is to be ontologically determinate in some respect. &
        Identity is not an attribute added to an independently constituted entity. Being and identity are inseparable in the actuality of an entity. \\

        Constitutive determination &
        Determination concerns the definite character through which an entity is what it is. It belongs to the ontological constitution of the entity rather than merely to its external description. &
        Identity is inseparable from the determinations constitutive of the entity, although it cannot be reduced to an inventory of separately listed determinations. \\

        Determinate character &
        An entity is not an entirely indeterminate possibility. Its actuality involves a definite ontological character, even when that character is not exhaustively understood. &
        Identity presupposes determinateness, but the complete articulation of an entity's identity need not be available or reducible to a finite description. \\

        Internal determination &
        Some determinations belong to the constitution of an entity itself rather than arising solely from its relations to other entities. &
        The identity of an entity cannot be explained exclusively through its external relations or contextual descriptions. \\

        Relational determination &
        An entity may acquire or exhibit determinate conditions through its relations with other entities and domains. &
        Relational determinations may affect the entity's condition and expression without necessarily replacing or exhausting its identity. \\

        Structural determination &
        An entity may be constituted or constrained by an organized configuration of determinations and their relations. &
        Structural organization can contribute to identity, but identity should not be equated with a structural description alone. \\

        Material determination &
        An entity may possess determinations associated with its material constitution and realization. &
        Material characteristics may be constitutive of an entity at a particular ontological domain, but identity cannot be universally reduced to material properties. \\

        Determination and distinction &
        Determination establishes definite differences between an entity and what it is not, without requiring that every difference be identity-defining. &
        Not every distinguishing feature is constitutive of identity. The ontological significance of a determination must be established rather than assumed. \\

        Determination and change &
        Determinations may persist, vary, emerge, or cease through ontological processes. &
        Changes in determination raise the question of identity preservation, but neither every change destroys identity nor every persistence guarantees it. \\

        Determination and identity preservation &
        The continuity of an entity may involve the preservation of some determinations and the transformation of others. &
        Identity preservation requires an account of which determinations are constitutive and under what conditions their transformation remains compatible with the entity's identity. \\

        Limits of determination &
        No isolated determination necessarily provides an exhaustive account of an entity's ontological constitution. &
        Identity cannot be reduced to a single property, feature, relation, structural description, or material condition without further argument. \\

        \hline
    \end{tabular}
\end{table}
